\section{The smooth case: groundwork for Theorem 3.1 (projective space)}\label{sec:smooth-blueprint}

Section~5 uses the envelope property of $\PP^n$ (FHX Theorem~3.1). For $n\ge1$, the numerical
and rational Picard groups identify with $\R$ and $\Q$, respectively. For $n=0$ both geometric
groups vanish; the scalar parameters in the model below remain formal parameters in that case.
The author decided to formalize
Section~3 before Proposition~\ref{prop:nonflat-norm}; the formalization starts with $X=\PP^n$, where
line bundles are $\mathcal O(a)$, sections are forms, and $S=X\times\A^1$ is computed in
$k[x_0,\dots,x_n,t]$. In characteristic $p$ the proof is FHX's test-ideal argument; in characteristic
$0$ it is the $L^2$ argument of Section~\ref{sec:char0-blueprint} (Ohsawa--Takegoshi and H\"ormander estimates on the affine
cone, from the \texttt{TranscendentalNakamaye} project), which replaces \cite[Theorem~5.20]{BJ22} on $\PP^n$.
General smooth projective varieties are treated in Section~\ref{sec:smooth-general}.

\begin{definition}[The analytification of $\PP^n$]\label{def:pn-an}
\lean{EnvelopeProperty.ProjBerk, EnvelopeProperty.ProjBerk.instCompactSpace, EnvelopeProperty.ProjBerk.triv}
\leanok
\source{fhx-envelope-20261028:subsec:prelim-berkovich}
$(\PP^n)^\an$ is modelled by the semivaluations $|\cdot|$ of $k[x_0,\dots,x_n]$ trivial on $k$ that are
graded ($|f|=\max_d|f_d|$) and normalized ($\max_i|x_i|=1$); it is compact Hausdorff. A point
$(\xi,|\cdot|)$ of the Berkovich analytification gives $|f|=|f/x_j^d|/\max_i|x_i/x_j|^d$ on forms of degree
$d$; conversely the kernel is a relevant homogeneous prime. (The comparison with the analytification of
the scheme $\PP^n$ is part of the general Berkovich groundwork, not yet formalized.)
\end{definition}

\begin{definition}[Gauss extension and $\phi_I$]\label{def:pn-gauss}
\lean{EnvelopeProperty.ProjBerk.gaussExt, EnvelopeProperty.ProjBerk.phiIdeal, EnvelopeProperty.ProjBerk.phiIdeal_mul}
\leanok
\uses{def:pn-an}
\source{fhx-envelope-20261028:subsec:prelim-berkovich}
$\sigma(v)(\sum_\lambda s_\lambda t^\lambda)=\max_\lambda|s_\lambda|_ve^{-\lambda}$, and for an ideal $I$ of
$k[x][t]$, $\phi_I(v)=\sup_{g\in I}\log|g|_{\sigma(v)}$; $\phi_{IJ}=\phi_I+\phi_J$.
\end{definition}

\begin{definition}[Fubini--Study and $\theta$-psh functions on $\PP^n$]\label{def:pn-psh}
\lean{EnvelopeProperty.ProjBerk.FSData, EnvelopeProperty.ProjBerk.IsFS, EnvelopeProperty.ProjBerk.IsRatFS, EnvelopeProperty.ProjBerk.IsPSH, EnvelopeProperty.ProjBerk.IsFS.isPSH, EnvelopeProperty.ProjBerk.isPSH_const}
\leanok
\uses{def:pn-an}
\source{fhx-envelope-20261028:subsec:prelim-psh, bj22-trivval:defi:psh}
FS functions $m^{-1}\max_j\{\log|s_j|+\lambda_j\}$ for $\mathcal O(a)$; $\theta$-psh functions
with scalar parameter $\theta\in\R$ (identified with $N^1(\PP^n)_\R$ when $n\ge1$):
generically finite usc decreasing limits of nets of FS functions with
classes tending to $\theta$ \cite[Definition~4.1]{BJ22}. FS functions are psh (PSH1); constants are
$\theta$-psh for $\theta\ge0$.
\end{definition}

\begin{definition}[Envelopes and the envelope property]\label{def:pn-ep}
\lean{EnvelopeProperty.ProjBerk.envelope, EnvelopeProperty.ProjBerk.HasEnvelopeProperty, EnvelopeProperty.uscReg, EnvelopeProperty.ProjBerk.IsRatPL}
\leanok
\uses{def:pn-psh}
\source{fhx-envelope-20261028:subsec:prelim-envelopes, bj22-trivval:sec:envprop}
$\mathrm P_\theta(\phi)=\sup\{\tau\in\PSH(\theta):\tau\le\phi\}$; $\theta$ has the envelope property if the
usc regularization of the supremum of every bounded-above family in $\PSH(\theta)$ is $\theta$-psh.
Rational PL functions are differences of finite-valued rational FS functions.
\end{definition}

\begin{lemma}[Frobenius roots]\label{lem:frobenius-roots}
\lean{EnvelopeProperty.frobRoot, EnvelopeProperty.frobTrace, EnvelopeProperty.gc_frobRoot_frobPow, EnvelopeProperty.frobRoot_mul_le, EnvelopeProperty.HasWeightedGeneratorsLE.frobRoot}
\leanok
\source{fhx-envelope-20261028:eq:test-def}
In a polynomial ring $R$ over a perfect field of characteristic $p$, $I^{[1/q]}=\Phi_e(I)$ for the
Frobenius trace $\Phi_e$ (trivialization $dx$ of $\omega_R$) is the smallest ideal $J$ with
$I\subseteq J^{[q]}$; $(IJ)^{[1/q]}\subseteq I^{[1/q]}J^{[1/q]}$; for weights $w$, an ideal generated in
$w$-degree $\le D$ has root generated in degree $\le\lfloor D/q\rfloor$.
\end{lemma}

\begin{lemma}[Test ideals and asymptotic test ideals]\label{lem:test-ideals}
\lean{EnvelopeProperty.testIdeal, EnvelopeProperty.exists_testIdeal_eq, EnvelopeProperty.testIdeal_add_le, EnvelopeProperty.asymptoticTestIdeal, EnvelopeProperty.le_asymptoticTestIdeal, EnvelopeProperty.asymptoticTestIdeal_mul_le, EnvelopeProperty.HasWeightedGeneratorsLE.asymptoticTestIdeal}
\leanok
\uses{lem:frobenius-roots}
\source{fhx-envelope-20261028:eq:smooth-test-inclusions}
$\tau(I^c)=\bigcup_e(I^{\lceil cp^e\rceil})^{[1/p^e]}$, attained for $e\gg0$; subadditivity
$\tau(I^{c+d})\subseteq\tau(I^c)\tau(I^d)$; for a graded sequence $J_\bullet$, $B_m=\tau(J_q^{m/q})$ for $q$
divisible, $J_m\subseteq B_m$, $B_{m\ell}\subseteq B_m^\ell$, and $B_m$ is generated in degree $\le mD$ when
$J_q$ is generated in degree $\le qD$ (the $\PP^n$ case of FHX Lemma~3.2, where $K_X+(n+1)H=0$).
\end{lemma}

\begin{lemma}[Divisorial trace estimate]\label{lem:pn-divisorial-trace}
\lean{EnvelopeProperty.Cartier.exists_logDiscrepancy_polynomial, EnvelopeProperty.Cartier.exists_logDiscrepancy, EnvelopeProperty.Cartier.IsPBasis.cartier_eq, EnvelopeProperty.Cartier.exists_isPBasis_of_essFiniteType, EnvelopeProperty.Cartier.ordGE_cartier, EnvelopeProperty.Cartier.le_scaled_of_ordGE}
\leanok
\uses{lem:test-ideals}
\source{fhx-envelope-20261028:lem:divisorial-trace-estimate}
For every discrete valuation ring $R'$ essentially of finite type over $k$, with
\[
 k[x_0,\dots,x_n,t]\subseteq R'\subseteq k(x,t),
\]
there is $A\ge0$ (the log discrepancy of $\ord_{R'}$) with
\[
 \ord_{R'}(\Phi_e f)\ge\ord_{R'}(f)/p^e-A
\]
for all $e$ and $f$ (FHX Lemma~3.3 for the divisorial points of
Definition~\ref{def:pn-div}). Formalized with the sharp constant:
\[
 \ord_{R'}(\Phi_ef)\ge\lfloor(\ord_{R'}f+a)/p^e\rfloor-a,
\]
where $a=\ord_{R'}(dx/du)$ for a $p$-basis $u$ of $R'$ containing a uniformizer,
so $A_S(\ord_{R'})=1+a$. The Cartier operator on top forms is independent of the $p$-basis.
\end{lemma}

\begin{theorem}[Envelope property of $\PP^n$ in characteristic $p$]\label{thm:pn-ep-charp}
\lean{EnvelopeProperty.ProjBerk.hasEnvelopeProperty_of_charP, EnvelopeProperty.ProjBerk.hasEnvelopeProperty_rat_of_charP, EnvelopeProperty.ProjBerk.isPSH_envelope_nat_of_nonpos, EnvelopeProperty.ProjBerk.isPSH_testLimit, EnvelopeProperty.ProjBerk.testLimit_le_envelope_of_isDivisorial, EnvelopeProperty.ProjBerk.exists_phiIdeal_asymptoticTestIdeal_le, EnvelopeProperty.ProjBerk.exists_le_ord_frobTrace, EnvelopeProperty.ProjBerk.hasEnvelopeProperty_of_neg}
\leanok
\uses{def:pn-ep, lem:test-ideals, lem:pn-divisorial-trace, rmk:bj-pn-ep-assembly}
\source{fhx-envelope-20261028:thm:smooth-EP}
Let $k$ be algebraically closed of characteristic $p>0$. Every $\theta\ge0$ has the envelope property on $\PP^n$.
\end{theorem}

\begin{proof}
Remark~\ref{rmk:bj-pn-ep-assembly}, steps (a)--(e). In (d) the bound $\operatorname{ord}_{R'}(B_m)\ge(m/q)\operatorname{ord}_{R'}(J_q)-A$ is
proved generator by generator from $B_m=\Phi_e(J_q^{\lceil mp^e/q\rceil})$ with $A=1+a$ (Lemma~\ref{lem:pn-divisorial-trace}),
the rounding being harmless since $\phi_{J_q}\le0$; it holds for every graded sequence of nonzero ideals. In (c),
$u\ge\mathrm Q\ge\inf\phi$ (a constant is a competitor). No countability of $k$ is needed. For $\theta<0$ the
statement holds vacuously, as $\PSH(\theta)=\emptyset$.
\end{proof}

\begin{theorem}[Envelope property of $\PP^n$]\label{thm:pn-ep}
\lean{EnvelopeProperty.ProjBerk.hasEnvelopeProperty_of_countable, EnvelopeProperty.ProjBerk.hasEnvelopeProperty_of_charP, EnvelopeProperty.ProjBerk.hasEnvelopeProperty_of_charZero}
\leanok
\uses{def:pn-ep, thm:pn-ep-charp, thm:c0-pn-ep}
\source{fhx-envelope-20261028:thm:smooth-EP}
Every $\theta\ge0$ has the envelope property on $\PP^n$ (in characteristic $0$ for countable $k$, the convention of
\cite[\S5]{FHX}; in characteristic $p$ for every algebraically closed $k$).
\end{theorem}
